\documentclass[11pt]{article}
\usepackage{graphicx}
\usepackage{amsmath}
\usepackage{hyperref}
\usepackage{booktabs}
\title{Physics-Informed Multi-Modal Foundation Models for Real-Time Solar-Wind Hybrid Forecasting}
\author{Authors}
\date{\today}
\begin{document}
\maketitle
\begin{abstract}
We propose a hybrid physics-informed, multi-modal foundation model that combines Physics-Informed Neural Networks (PINNs) with Vision Transformer (ViT) backbones and multi-level fusion transformers for real-time solar and wind hybrid forecasting. Key contributions: (1) hybrid loss combining data-driven and physics constraints, (2) temporal-spatial attention for multi-site forecasting, and (3) cross-domain pre-training for improved generalization.
\end{abstract}

\section{Introduction}
Joint solar and wind forecasting improves grid balancing by leveraging cross-modal information from gridded meteorological fields and local station measurements. Physics-informed losses ensure physically plausible predictions while allowing data-driven flexibility.

\section{Method}
Our prototype uses a ViT-like patch encoder for ERA5-like gridded inputs, a scalar MLP for local station measurements (irradiance and wind speed), and a fusion head producing joint solar and wind power forecasts. Physics losses include a solar irradiance upper bound and a turbine power-curve regularization term.

\section{Experiments}
We evaluate a small prototype on synthetic data designed to mimic irradiance and wind-speed driven generation. The training script emits scatter plots and attention visualizations for model introspection.

\section{Results}
Quantitative metrics (synthetic run):
\begin{table}[h]
\centering
\begin{tabular}{lrr}
\toprule
Metric & Solar (W) & Wind (W) \\
\midrule
RMSE & 6.35 & 751.13 \\
MAE  & 5.36 & 641.68 \\
\bottomrule
\end{tabular}
\caption{Prototype performance on synthetic test split.}
\end{table}

Figure~\ref{fig:pred_vs_gt} shows predicted vs ground truth scatter plots for solar and wind power. Figures~\ref{fig:attn0}--\ref{fig:attn3} show per-layer attention overlays.

\begin{figure}[ht]
\centering
\includegraphics[width=0.9\textwidth]{../outputs/pred_vs_gt.png}
\caption{Predicted vs ground-truth (solar left, wind right).}
\label{fig:pred_vs_gt}
\end{figure}

\begin{figure}[ht]
\centering
\includegraphics[width=0.45\textwidth]{../outputs/attn_layer_0_head_0.png}
\caption{Attention overlay: layer 0, head 0.}
\label{fig:attn0}
\end{figure}

\begin{figure}[ht]
\centering
\includegraphics[width=0.45\textwidth]{../outputs/attn_layer_1_head_0.png}
\caption{Attention overlay: layer 1, head 0.}
\label{fig:attn1}
\end{figure}

\begin{figure}[ht]
\centering
\includegraphics[width=0.45\textwidth]{../outputs/attn_layer_2_head_0.png}
\caption{Attention overlay: layer 2, head 0.}
\label{fig:attn2}
\end{figure}

\begin{figure}[ht]
\centering
\includegraphics[width=0.45\textwidth]{../outputs/attn_layer_3_head_0.png}
\caption{Attention overlay: layer 3, head 0.}
\label{fig:attn3}
\end{figure}

\section{Conclusion}
This prototype demonstrates how physics-informed losses and multi-modal fusion can produce reasonable joint forecasts with interpretable attention maps. Future work will replace synthetic data with ERA5/NSRDB/PV-Live/Wind DB and scale the model for production.

\appendix
\section*{Appendix: Implementation Details}

This appendix lists implementation choices, hyperparameters, and file locations to reproduce experiments.

\begin{itemize}
  \item Framework: PyTorch
  \item ViT encoder: patch size 8, embedding dim 128, 4 transformer layers, 4 attention heads
  \item Scalar encoder: 2-layer MLP, hidden 64
  \item Fusion head: concatenation followed by 2-layer MLP
  \item Physics losses: solar bounds and wind power-curve regularization (see `src/model.py`)
  \item Synthetic dataset: `src/data.py` generates images and targets; real-data ingestion stubs in `src/datasets_real.py`
  \item Training script: `src/train.py` (use `--config` to point to `data_config_example.yaml` for real data)
  \item Outputs: `outputs/` contains trained weights and figures
\end{itemize}

Hyperparameters used for the prototype run
\begin{itemize}
  \item Learning rate: 1e-3
  \item Batch size: 32
  \item Epochs: 5
  \item Loss weights: see inline comments in `src/train.py`
\end{itemize}


\end{document}
